\section{سؤالی با چند بخش (۱۲ نمره)}

صورت سؤال اینجا می‌آید. هر \lr{\texttt{section}} یک سؤال است: شماره در کادری
تیره می‌نشیند، عنوان کنارش می‌آید و خطی سایه‌دار زیر آن کشیده می‌شود.

\begin{enumerate}
    \item زیرسؤال اول.
    \item زیرسؤال دوم، همراه با پانوشت برای معادل
          لاتین.\footnote{\lr{Footnote}}
    \item زیرسؤال سوم.
\end{enumerate}

\section{سؤالی با زیربخش (۱۰ نمره)}

وقتی یک سؤال به بخش‌هایی تقسیم می‌شود که هرکدام عنوان خود را می‌خواهند، از
زیربخش استفاده کنید.

\subsection{بخش نخست}

شماره‌ی زیربخش‌ها در کادری توخالی می‌آید، نه پُر، تا دو سطح از هم قابل تشخیص
بمانند.

\subsection{بخش دوم}

متن با قلم \lr{XM Yekan} و عنوان‌ها و شماره‌ها با قلم پیکسلی سورنا چیده
می‌شوند.

\section{سؤالی با کادر راهنما (۲۰ نمره)}

\begin{neobox}[neotitlefont=\pixelfont]{راهنمایی}
    \lr{\texttt{neobox}} همان کادری است که راهنمای ابتدای تمرین را نگه می‌دارد.
    با کلید \lr{\texttt{neotitlefont}} قلم عنوان و با آرگومان اختیاری هر کلید
    \lr{\texttt{tcolorbox}} را می‌توانید بدهید.
\end{neobox}
